% Contenidos del capítulo.
% Las secciones presentadas son orientativas y no representan
% necesariamente la organización que debe tener este capítulo.

\section{Introducción}

El presente trabajo de fin de grado describe el desarrollo de una \textbf{tienda online} para la \textbf{Droguería Mercería Inés}, un negocio familiar ubicado en Valencia. El establecimiento, regentado por Ester, ofrece una amplia variedad de productos: perfumes, utensilios de costura, hilos, telas especiales y artículos de mercería en general, siguiendo la tradición de las droguerías de barrio que durante décadas han formado parte del comercio local.

Aunque el negocio cuenta con presencia en redes sociales, no disponía de una página web propia que permitiese a sus clientes consultar el catálogo, realizar pedidos y efectuar pagos de forma online. Este proyecto nace con el objetivo de cubrir esa necesidad mediante el desarrollo de una aplicación web completa, construida íntegramente con herramientas gratuitas y de código abierto.

La solución desarrollada permite a los clientes registrarse, explorar el catálogo de productos, gestionar su carrito y completar compras a través de la pasarela de pago \textit{Stripe}. Por parte de la administración, el sistema facilita la gestión del inventario y el seguimiento de pedidos, con notificaciones automáticas de confirmación de compra y avisos de bajo stock a través de un \textit{bot de Telegram}.

\section{Motivación}

La principal motivación de este proyecto es de carácter personal: dar respuesta a una necesidad real de un negocio familiar. La Droguería Mercería Inés tiene presencia en redes sociales, pero carecía de una plataforma web propia desde la que sus clientes pudieran realizar compras de manera cómoda y segura. Desarrollar esta solución supone una aportación directa al negocio de la familia, facilitando tanto la experiencia de compra de los clientes como el trabajo diario de la propietaria en la gestión del stock y los pedidos.

Además, existe una motivación de carácter profesional. La elección de \textbf{Angular} y \textbf{Spring Boot} como tecnologías principales responde al hecho de que son las herramientas utilizadas en \textbf{Capgemini}, empresa en la que se van a realizar las prácticas. Abordar este proyecto ha permitido adquirir una base sólida en ambas tecnologías antes de incorporarse a un entorno profesional, complementando la formación académica con una experiencia de desarrollo real y completa.

Por último, este proyecto pretende demostrar que es posible construir una solución de comercio electrónico funcional y profesional haciendo uso exclusivamente de herramientas \textbf{gratuitas y de código abierto}, sin necesidad de recurrir a plataformas de pago ni a infraestructuras costosas.

\section{Objetivos}

El objetivo principal del proyecto es el siguiente:

\begin{itemize}
    \item \textbf{Desarrollar una tienda online funcional y completa para la Droguería Mercería Inés}, que permita a sus clientes consultar el catálogo, realizar pedidos y efectuar pagos de forma segura a través de internet.
\end{itemize}

Para alcanzarlo, se han definido los siguientes objetivos secundarios:

\begin{itemize}
    \item Dar visibilidad online al negocio mediante una plataforma web propia, complementando su presencia actual en redes sociales.
    \item Facilitar a la propietaria el control del inventario y el seguimiento de los pedidos recibidos a través de la aplicación.
    \item Automatizar las notificaciones de confirmación de pedido y los avisos de bajo stock mediante un bot de Telegram integrado en el sistema.
    \item Adquirir experiencia práctica en \textbf{Angular} y \textbf{Spring Boot} como preparación para las prácticas profesionales en Capgemini.
    \item Demostrar que una solución de comercio electrónico completa puede desarrollarse íntegramente con herramientas gratuitas y de código abierto.
\end{itemize}

\section{Organización de la memoria}

La memoria se organiza en los siguientes capítulos:

\begin{itemize}
    \item \textbf{Capítulo 1 — Introducción:} presenta el proyecto, su motivación, los objetivos planteados y la estructura del documento.
    \item \textbf{Capítulo 2 — Estado del arte:} analiza las plataformas de creación de tiendas online más representativas y estudia las tecnologías seleccionadas para el desarrollo del sistema, justificando cada elección.
    \item \textbf{Capítulo 3 — Requisitos, especificaciones, coste, riesgos y viabilidad:} recoge los requisitos funcionales y no funcionales del sistema, las especificaciones técnicas, la planificación temporal y económica del proyecto, los riesgos identificados y el análisis de viabilidad.
    \item \textbf{Capítulo 4 — Análisis:} incluye los diagramas de análisis del sistema: el diagrama de casos de uso, el diagrama de clases de análisis y los diagramas de secuencia general del sistema (DSGS).
    \item \textbf{Capítulo 5 — Diseño:} describe la arquitectura general del sistema, los diagramas UML de diseño, el modelo de la base de datos y los prototipos de la interfaz de usuario.
    \item \textbf{Capítulo 6 — Implementación y pruebas:} detalla el proceso de desarrollo de la aplicación y los resultados de las pruebas funcionales, de rendimiento y de usabilidad realizadas.
    \item \textbf{Capítulo 7 — Conclusiones:} revisa los costes finales del proyecto, presenta las conclusiones obtenidas y propone posibles líneas de trabajo futuro.
\end{itemize}
